\chapter{Control Design and Analysis}
\label{ch:control}

\section{The design problem}

Chapter~\ref{ch:converter} established the difficulty precisely. The plant
$G_{vd}$ is a second-order resonance whose frequency is essentially fixed at
about \SI{412}{\hertz} but whose damping ratio varies from \num{0.023} to
\num{1.035} depending only on where the array sits on its I--V curve. A
controller must hold adequate stability margins across that entire range, using
one fixed set of gains, because the operating point is dictated by the MPPT
layer and by the weather rather than by the designer.

Two candidate architectures are compared.

\begin{figure}[htbp]
\centering
\resizebox{\textwidth}{!}{%
\begin{tikzpicture}[
  scale=0.88, transform shape,
  font=\small,
  box/.style={draw, rounded corners=2pt, minimum height=9mm, minimum width=24mm, align=center},
  >={Stealth[length=2.2mm]}
]
\node[box] (mppt) {MPPT\\ \SI{50}{\hertz}};
\node[box, right=10mm of mppt] (vloop) {Voltage PI\\ \SI{250}{\hertz}};
\node[box, right=10mm of vloop] (iloop) {Current PI\\ \SI{2}{\kilo\hertz}};
\node[box, right=10mm of iloop] (pwm) {PWM\\ \SI{20}{\kilo\hertz}};
\node[box, right=10mm of pwm] (plant) {Converter\\ and array};

\draw[->] (mppt) -- node[above]{$V^\ast$} (vloop);
\draw[->] (vloop) -- node[above]{$i_L^\ast$} (iloop);
\draw[->] (iloop) -- node[above]{$d$} (pwm);
\draw[->] (pwm) -- (plant);

\coordinate (out) at ($(plant.east)+(9mm,0)$);
\draw[->] (plant.east) -- (out) node[right]{$\vpv,\ \ipv,\ i_L$};

\draw[->] ($(plant.east)+(5mm,0)$) |- ($(iloop.south)+(0,-7mm)$) -| (iloop.south)
  node[pos=0.25, below]{$i_L$};
\draw[->] ($(plant.east)+(7mm,0)$) |- ($(vloop.south)+(0,-14mm)$) -| (vloop.south)
  node[pos=0.25, below]{$\vpv$};
\draw[->] ($(plant.east)+(9mm,0)$) |- ($(mppt.south)+(0,-21mm)$) -| (mppt.south)
  node[pos=0.25, below]{$\vpv,\ \ipv$};
\end{tikzpicture}}
\caption[Cascaded control architecture]{The cascaded architecture. Each layer commands the reference of the layer to its left and runs roughly eight times slower than the layer to its right.}
\label{fig:cascade}
\end{figure}

\begin{description}[leftmargin=0pt, style=unboxed]
\item[Single-loop.] One proportional--integral controller acting on the PV
  voltage error and producing the duty ratio directly. This is the obvious first
  attempt and appears widely in the MPPT literature.
\item[Cascade.] An inner PI regulating inductor current, with an outer PI
  regulating PV voltage by commanding the inner current reference, as in
  Figure~\ref{fig:cascade}.
\end{description}

\section{The single-loop design and why it fails}

\subsection{Design}

The single-loop controller was designed conventionally: crossover placed at
\SI{130}{\hertz}, roughly a third of the LC resonance, with the PI zero a factor
of five below crossover at \SI{26}{\hertz}. Evaluating the plant magnitude at
crossover and solving $|L(j\omega_c)| = 1$ gives $K_p = \num{2.264e-3}$ and
$K_i = \num{0.370}$. Designed at the STC maximum power point, this is a
perfectly reasonable-looking design.

\subsection{The multiple-crossover problem}

The trouble appears immediately when the loop gain is examined over frequency
rather than at a single point. Because a PI contributes an approximately flat
gain above its zero, and because the plant magnitude is flat below its
resonance, the loop gain is nearly flat between the intended crossover and the
resonance. The resonance then lifts the magnitude by up to \SI{27}{\decibel},
pushing $|L|$ back above unity, and it must come down through unity a third
time on the far side, by which point the plant has contributed its full
\SI{-180}{\degree}.

The measured result is three gain crossovers, at \SI{56}{\hertz},
\SI{130}{\hertz} and \SI{542}{\hertz}, with a peak loop gain of
\SI{43.4}{\decibel}. When a loop crosses unity more than once, quoting the
margin at the intended crossover is meaningless; the relevant number is the
worst margin over all crossings.

\subsection{Collapse of the phase margin}

Table~\ref{tab:singleloop} gives the worst-case phase margin along the I--V
curve. On the current-source branch, where the damping ratio is \num{0.023}, the
phase margin falls to \SI{1.5}{\degree}. The loop is nominally stable, but a
\SI{1.5}{\degree} margin is no margin at all: any unmodelled delay, any sensor
filter, any component tolerance destroys it, and even nominally it would ring
almost undamped for a very long time.

\begin{table}[htbp]
\centering
\caption{Worst-case phase margin along the I--V curve at \SI{1000}{\watt\per\square\metre}. The single-loop controller degrades exactly where the array behaves as a current source.}
\label{tab:singleloop}
\begin{tabular}{@{}rrrrrr@{}}
\toprule
$V$ (\si{\volt}) & $r_d$ (\si{\ohm}) & $\zeta$ & \multicolumn{2}{c}{Single-loop PI} & Cascade outer \\
\cmidrule(lr){4-5}\cmidrule(lr){6-6}
 & & & PM (\si{\degree}) & crossings & PM (\si{\degree}) \\
\midrule
69.09  & 482.1 & 0.0234 & 1.46  & 3 & 85.3 \\
88.83  & 475.8 & 0.0234 & 1.47  & 3 & 85.3 \\
108.57 & 417.6 & 0.0240 & 1.57  & 3 & 85.4 \\
128.31 & 182.9 & 0.0299 & 2.61  & 3 & 86.4 \\
138.18 & 77.23 & 0.0444 & 5.14  & 3 & 88.8 \\
148.05 & 27.94 & 0.0884 & 12.91 & 3 & 95.9 \\
157.92 & 10.25 & 0.2068 & 34.50 & 3 & 114.2 \\
167.79 & 4.458 & 0.4463 & 88.62 & 3 & 122.7 \\
177.66 & 2.539 & 0.7600 & 139.8 & 1 & 110.6 \\
187.53 & 1.832 & 1.0350 & 135.4 & 1 & 105.3 \\
\bottomrule
\end{tabular}
\end{table}

This is not a matter of retuning. The problem is structural: a PI has no
mechanism for damping a resonance it cannot see, and the required gain reduction
to keep $|L|$ below unity through the resonance would push crossover down to a
few hertz, making the loop uselessly slow.

\section{The cascaded design}

\subsection{Why an inner current loop fixes the problem}

Consider the transfer function from duty ratio to inductor current,

\begin{equation}
G_{id}(s) = \frac{i_L(s)}{d(s)}
= \frac{V_{dc}}{L}\cdot
  \frac{s - g/C_{\mathrm{pv}}}
       {s^2 + \left(\tfrac{R_L}{L} - \tfrac{g}{C_{\mathrm{pv}}}\right)s + \tfrac{1-gR_L}{L C_{\mathrm{pv}}}} .
\label{eq:Gid}
\end{equation}

The zero sits at $s = g/C_{\mathrm{pv}}$, which is in the left half plane
because $g < 0$, so this path is minimum phase. More importantly, at frequencies
well above the resonance the expression collapses to

\begin{equation}
G_{id}(s) \;\approx\; \frac{V_{dc}}{L s},
\label{eq:Gidhf}
\end{equation}

a pure integrator whose gain depends only on $V_{dc}$ and $L$ and \emph{not at
all} on the operating point. This is visible directly in the bottom-left panel
of Figure~\ref{fig:plantbode}, where the six curves are indistinguishable above
a few hundred hertz.

So an inner loop closed at \SI{2}{\kilo\hertz}, comfortably above the
\SI{412}{\hertz} resonance, sees an operating-point-independent plant. Its gains
follow immediately from~\eqref{eq:Gidhf}: for unity loop gain at
$\omega_{ci} = 2\pi \times \SI{2000}{\hertz}$,

\begin{equation}
K_{p,i} = \frac{L\,\omega_{ci}}{V_{dc}} = \num{0.04712},
\qquad
K_{i,i} = K_{p,i}\frac{\omega_{ci}}{10} = \num{59.2},
\end{equation}

placing the PI zero a decade below crossover. Because the plant contributes
\SI{-90}{\degree} there and the PI adds only a few degrees more, the phase
margin is close to \SI{85}{\degree} by construction.

\subsection{The outer loop sees a first-order plant}

With the inner loop fast and accurate, $i_L \approx i_L^\ast$, and the input
capacitor equation~\eqref{eq:av1} linearises to

\begin{equation}
\frac{\tilde v_{\mathrm{pv}}(s)}{\tilde i_L^\ast(s)}
= \frac{-1}{C_{\mathrm{pv}} s - g}
= \frac{-1}{C_{\mathrm{pv}} s + 1/r_d} .
\label{eq:outerplant}
\end{equation}

The resonance is gone. What remains is a single real pole at
$1/(r_d C_{\mathrm{pv}})$. Its DC gain is $r_d$, which still varies widely, but
at the outer crossover of \SI{250}{\hertz} the capacitive term dominates
whenever $\omega \gg 1/(r_d C_{\mathrm{pv}})$, so the plant looks like
$1/(C_{\mathrm{pv}} s)$ and the required proportional gain is again
operating-point independent:

\begin{equation}
K_{p,v} = C_{\mathrm{pv}}\,\omega_{cv} = \num{0.15708},
\qquad
K_{i,v} = K_{p,v}\frac{\omega_{cv}}{5} = \num{49.35}.
\end{equation}

Physically, the inner current loop turns the inductor into a controlled current
source, and a controlled current source charging a capacitor is about the
easiest plant in control engineering. That is the whole reason the cascade
works.

\subsection{Verification on the full-order model}

The gains above were derived from simplified models: the high-frequency
approximation~\eqref{eq:Gidhf} for the inner loop and the ideal-inner-loop
approximation~\eqref{eq:outerplant} for the outer. Design on a simplified model
is only legitimate if the result is then checked on the full one.

Both loops were therefore re-analysed on the exact second-order plant, with the
inner loop closed exactly rather than assumed ideal:

\begin{align}
L_i(s) &= C_i(s)\,G_{id}(s), \\
L_v(s) &= -C_v(s)\cdot G_{vd}(s)\,\frac{C_i(s)}{1 + C_i(s) G_{id}(s)} .
\end{align}

Results are in the last column of Table~\ref{tab:singleloop} and in
Table~\ref{tab:cascadeG}. The inner loop margin is between \SI{85.0}{\degree}
and \SI{85.8}{\degree} everywhere, confirming operating-point independence. The
outer loop margin is between \SI{85.3}{\degree} and \SI{122.7}{\degree}, always
far above the \SI{45}{\degree} target.

\begin{table}[htbp]
\centering
\caption{Cascade margins at the maximum power point across irradiance, verified on the full-order plant.}
\label{tab:cascadeG}
\begin{tabular}{@{}rrrrrr@{}}
\toprule
$G$ (\si{\watt\per\square\metre}) & $r_d$ (\si{\ohm}) & Inner PM (\si{\degree}) & Inner $f_c$ (\si{\hertz}) & Outer PM (\si{\degree}) & Outer $f_c$ (\si{\hertz}) \\
\midrule
200  & 50.97 & 85.02 & 2090 & 90.9  & 215.7 \\
400  & 25.95 & 85.06 & 2090 & 96.8  & 205.0 \\
600  & 17.38 & 85.09 & 2090 & 102.6 & 191.4 \\
800  & 13.02 & 85.12 & 2090 & 108.4 & 174.7 \\
1000 & 10.37 & 85.16 & 2089 & 113.9 & 155.4 \\
\bottomrule
\end{tabular}
\end{table}

\begin{figure}[htbp]
\centering
\bigfig{s03_control_design.png}
\caption[Control design comparison]{Comparison of the two architectures. Top left, single-loop gain magnitude at five operating points: the resonance forces the curve back above \SI{0}{\decibel}. Top right, the cascade outer loop crosses unity once and cleanly. Bottom left, worst-case phase margin against operating voltage, the central result of the chapter. Bottom right, closed-loop reference step responses of the cascade at five operating points.}
\label{fig:controldesign}
\end{figure}

\subsection{Bandwidth variation and what it costs}

One genuine operating-point dependence remains in the cascade: the outer
crossover falls from about \SI{223}{\hertz} on the current-source branch to
\SI{15}{\hertz} near $\Voc$. This happens because the plant DC gain $r_d$ has
become so small there that the loop gain drops below unity at a lower frequency.
The bottom-right panel of Figure~\ref{fig:controldesign} shows the
consequence directly: the settling time varies from a few milliseconds to about
\SI{50}{\milli\second} across the range.

This variation is benign, because the slow region is far from the maximum power
point where the tracker actually operates, and because the phase margin remains
above \SI{100}{\degree} throughout. If uniform bandwidth were required, the
outer proportional gain could be scheduled on the measured $r_d$, which is
directly available as $-\Delta \vpv / \Delta \ipv$ from the MPPT perturbations.
That refinement was not needed here and is noted as future work.

\subsection{Measured step response}

The design predicts a \SI{250}{\hertz} outer loop with a large phase margin,
therefore a well-damped response settling in roughly \SI{10}{\milli\second}.
Simulating the full nonlinear averaged plant with the discrete controllers, a
\SI{140}{\volt} to \SI{165}{\volt} reference step settles to within
\SI{2}{\percent} in \SI{11.9}{\milli\second} with no overshoot, and the
steady-state error is zero to solver precision. Measured against the analytic
I--V curve, the settled operating point matches to better than
\SI{0.05}{\percent} in both voltage and current.

\section{Consequences for the MPPT layer}

The voltage loop settling time of roughly \SI{12}{\milli\second} is not merely a
performance figure; it is a hard constraint on the layer above. A perturb and
observe tracker must wait for the voltage loop to settle before it measures the
power resulting from its perturbation, otherwise it measures a transient and
draws the wrong conclusion. That is why the MPPT sample period was set to
\SI{20}{\milli\second}, comfortably longer than the settling time.

This is timescale separation appearing as a design rule rather than as an
abstraction, and it is the reason a control-theoretic treatment of MPPT is
worth the effort: the number \SI{20}{\milli\second} is derived, not guessed.
